% =====================================================================
%  conclusion.tex  --  PHASE VII, Task 14C
%
%  Section 12 (\label{sec:conclusion}).  Written after the limitations
%  section, so it does not have to hedge: the boundaries are already
%  stated and the conclusion can be direct about what was and was not
%  demonstrated.
%
%  Drop-in usage: % =====================================================================
%  conclusion.tex  --  PHASE VII, Task 14C
%
%  Section 12 (\label{sec:conclusion}).  Written after the limitations
%  section, so it does not have to hedge: the boundaries are already
%  stated and the conclusion can be direct about what was and was not
%  demonstrated.
%
%  Drop-in usage: % =====================================================================
%  conclusion.tex  --  PHASE VII, Task 14C
%
%  Section 12 (\label{sec:conclusion}).  Written after the limitations
%  section, so it does not have to hedge: the boundaries are already
%  stated and the conclusion can be direct about what was and was not
%  demonstrated.
%
%  Drop-in usage: % =====================================================================
%  conclusion.tex  --  PHASE VII, Task 14C
%
%  Section 12 (\label{sec:conclusion}).  Written after the limitations
%  section, so it does not have to hedge: the boundaries are already
%  stated and the conclusion can be direct about what was and was not
%  demonstrated.
%
%  Drop-in usage: \input{conclusion}
%  Bibliography: paper/references.bib
% =====================================================================

\section{Conclusion}
\label{sec:conclusion}

This paper set out to connect two capabilities that had matured
separately: inferring an industrial machine's operating state from its
electrical measurements, and deciding optimally when to de-energise it.
The connecting quantity is the duration of a non-productive interval,
which the first literature does not forecast and the second assumes is
given. We proposed a unified label-free framework that supplies it, and
evaluated every stage end to end on a public industrial benchmark of
eight machines.

Four results stand.

\textbf{Label-free state inference can be made trustworthy, but only if
the fit is gated.} Fitted at ten seeds on identical data, the labeller
places the standby cluster on a physically different state in three of
them, and the headline standby total moves by a third of its value.
Five circuit-level checks --- none of which uses the model's own
statistics --- separate those three from the other seven exactly, and
conditioning on them collapses the spread to $3\%$. An accept/reject
step is therefore part of the method, not an appendix to it, and this is
the paper's central methodological claim.

\textbf{Temporal coherence requires a duration model, not a stronger
prior.} A transition matrix can penalise a state change by at most
$\log(p_{\text{self}}/p_{\text{switch}})$ --- $9.23$ nats here --- while
the emission term routinely exceeds that by orders of magnitude, so at
over half of all timesteps the transition model cannot influence the
decode at all. Viterbi decoding takes flicker from $66.9\%$ to $44.5\%$;
an explicit minimum-dwell constraint takes it to zero, at a cost of
$0.19\%$ of rows and no measurable change to any physics or schedule
check, on all eight machines.

\textbf{Two of the framework's own components did not earn their
place, and are reported as such.} An untrained persistence reference
beats all five neural sequence forecasters, the originally proposed
sequence-to-sequence model by the largest margin of the seven; only a
gradient-boosted tree clears the reference, and by $1.45$\,s of
per-window MAE. On the economics, the proposed policy cannot be shown to
beat the plant's own behaviour on an 18-day held-out sample. What it
does beat, significantly and in every scenario, is the fixed-threshold
rule it was built to replace. The defensible decision-layer claim is
that \emph{adaptive thresholds beat a fixed one}, not that forecasting
beats not forecasting.

\textbf{The framework says where it does not apply.} Across the plant
the state definition is physically valid on the four motor-driven
machines and is rejected by the battery on the four auxiliary ones,
where no energised-idle state exists to find. Applied to a second site
it correctly characterises a CNC machining centre and --- on a
single-channel record, having lost the two checks that carry the
discriminating power --- fails to reject a solar inverter, a device that
cannot be shut down at all. A validation battery is only worth quoting
if it can reject something; this one rejects four machines out of eight
and, under reduced instrumentation, is shown to have a blind spot.

\medskip
\noindent
Set against the improvement the framework can extract, the honest
summary is that the facility's total attainable standby saving is under
$90$\,USD/yr, because the plant already de-energises for its long idle
gaps. \textbf{On this facility the framework's value is validated state
inference and decision quality, not recovered energy}, and that
conclusion belongs in the abstract rather than at the end of a
discussion. Whether a plant that does not already shut down manually,
or one with a larger energised-idle fraction, changes the arithmetic is
exactly the question the method now makes answerable from a meter
alone.

\subsection*{Future work}

Four extensions follow directly from the limitations of
Section~\ref{sec:limitations}. A \emph{longer record} is the only thing
that can settle the economic question: eighteen factory days is the
binding constraint on every currency interval here, and no estimator
narrows it. A \emph{second six-channel industrial dataset} would test
the physics battery rather than only the pipeline, which the available
single-channel record cannot do. A \emph{hidden semi-Markov
formulation}~\cite{yu2010hsmm} would learn the state-duration
distribution instead of imposing a constant, and a per-machine state
count selected under the physics veto would remove the $k=4$ assumption
that the eight-machine study shows to be wrong in both directions.
Finally, a \emph{seed-resampled ablation of the validation battery
itself}, check by check, would price the component this paper argues is
most important and is the one component its ablation table does not
contain.

%  Bibliography: paper/references.bib
% =====================================================================

\section{Conclusion}
\label{sec:conclusion}

This paper set out to connect two capabilities that had matured
separately: inferring an industrial machine's operating state from its
electrical measurements, and deciding optimally when to de-energise it.
The connecting quantity is the duration of a non-productive interval,
which the first literature does not forecast and the second assumes is
given. We proposed a unified label-free framework that supplies it, and
evaluated every stage end to end on a public industrial benchmark of
eight machines.

Four results stand.

\textbf{Label-free state inference can be made trustworthy, but only if
the fit is gated.} Fitted at ten seeds on identical data, the labeller
places the standby cluster on a physically different state in three of
them, and the headline standby total moves by a third of its value.
Five circuit-level checks --- none of which uses the model's own
statistics --- separate those three from the other seven exactly, and
conditioning on them collapses the spread to $3\%$. An accept/reject
step is therefore part of the method, not an appendix to it, and this is
the paper's central methodological claim.

\textbf{Temporal coherence requires a duration model, not a stronger
prior.} A transition matrix can penalise a state change by at most
$\log(p_{\text{self}}/p_{\text{switch}})$ --- $9.23$ nats here --- while
the emission term routinely exceeds that by orders of magnitude, so at
over half of all timesteps the transition model cannot influence the
decode at all. Viterbi decoding takes flicker from $66.9\%$ to $44.5\%$;
an explicit minimum-dwell constraint takes it to zero, at a cost of
$0.19\%$ of rows and no measurable change to any physics or schedule
check, on all eight machines.

\textbf{Two of the framework's own components did not earn their
place, and are reported as such.} An untrained persistence reference
beats all five neural sequence forecasters, the originally proposed
sequence-to-sequence model by the largest margin of the seven; only a
gradient-boosted tree clears the reference, and by $1.45$\,s of
per-window MAE. On the economics, the proposed policy cannot be shown to
beat the plant's own behaviour on an 18-day held-out sample. What it
does beat, significantly and in every scenario, is the fixed-threshold
rule it was built to replace. The defensible decision-layer claim is
that \emph{adaptive thresholds beat a fixed one}, not that forecasting
beats not forecasting.

\textbf{The framework says where it does not apply.} Across the plant
the state definition is physically valid on the four motor-driven
machines and is rejected by the battery on the four auxiliary ones,
where no energised-idle state exists to find. Applied to a second site
it correctly characterises a CNC machining centre and --- on a
single-channel record, having lost the two checks that carry the
discriminating power --- fails to reject a solar inverter, a device that
cannot be shut down at all. A validation battery is only worth quoting
if it can reject something; this one rejects four machines out of eight
and, under reduced instrumentation, is shown to have a blind spot.

\medskip
\noindent
Set against the improvement the framework can extract, the honest
summary is that the facility's total attainable standby saving is under
$90$\,USD/yr, because the plant already de-energises for its long idle
gaps. \textbf{On this facility the framework's value is validated state
inference and decision quality, not recovered energy}, and that
conclusion belongs in the abstract rather than at the end of a
discussion. Whether a plant that does not already shut down manually,
or one with a larger energised-idle fraction, changes the arithmetic is
exactly the question the method now makes answerable from a meter
alone.

\subsection*{Future work}

Four extensions follow directly from the limitations of
Section~\ref{sec:limitations}. A \emph{longer record} is the only thing
that can settle the economic question: eighteen factory days is the
binding constraint on every currency interval here, and no estimator
narrows it. A \emph{second six-channel industrial dataset} would test
the physics battery rather than only the pipeline, which the available
single-channel record cannot do. A \emph{hidden semi-Markov
formulation}~\cite{yu2010hsmm} would learn the state-duration
distribution instead of imposing a constant, and a per-machine state
count selected under the physics veto would remove the $k=4$ assumption
that the eight-machine study shows to be wrong in both directions.
Finally, a \emph{seed-resampled ablation of the validation battery
itself}, check by check, would price the component this paper argues is
most important and is the one component its ablation table does not
contain.

%  Bibliography: paper/references.bib
% =====================================================================

\section{Conclusion}
\label{sec:conclusion}

This paper set out to connect two capabilities that had matured
separately: inferring an industrial machine's operating state from its
electrical measurements, and deciding optimally when to de-energise it.
The connecting quantity is the duration of a non-productive interval,
which the first literature does not forecast and the second assumes is
given. We proposed a unified label-free framework that supplies it, and
evaluated every stage end to end on a public industrial benchmark of
eight machines.

Four results stand.

\textbf{Label-free state inference can be made trustworthy, but only if
the fit is gated.} Fitted at ten seeds on identical data, the labeller
places the standby cluster on a physically different state in three of
them, and the headline standby total moves by a third of its value.
Five circuit-level checks --- none of which uses the model's own
statistics --- separate those three from the other seven exactly, and
conditioning on them collapses the spread to $3\%$. An accept/reject
step is therefore part of the method, not an appendix to it, and this is
the paper's central methodological claim.

\textbf{Temporal coherence requires a duration model, not a stronger
prior.} A transition matrix can penalise a state change by at most
$\log(p_{\text{self}}/p_{\text{switch}})$ --- $9.23$ nats here --- while
the emission term routinely exceeds that by orders of magnitude, so at
over half of all timesteps the transition model cannot influence the
decode at all. Viterbi decoding takes flicker from $66.9\%$ to $44.5\%$;
an explicit minimum-dwell constraint takes it to zero, at a cost of
$0.19\%$ of rows and no measurable change to any physics or schedule
check, on all eight machines.

\textbf{Two of the framework's own components did not earn their
place, and are reported as such.} An untrained persistence reference
beats all five neural sequence forecasters, the originally proposed
sequence-to-sequence model by the largest margin of the seven; only a
gradient-boosted tree clears the reference, and by $1.45$\,s of
per-window MAE. On the economics, the proposed policy cannot be shown to
beat the plant's own behaviour on an 18-day held-out sample. What it
does beat, significantly and in every scenario, is the fixed-threshold
rule it was built to replace. The defensible decision-layer claim is
that \emph{adaptive thresholds beat a fixed one}, not that forecasting
beats not forecasting.

\textbf{The framework says where it does not apply.} Across the plant
the state definition is physically valid on the four motor-driven
machines and is rejected by the battery on the four auxiliary ones,
where no energised-idle state exists to find. Applied to a second site
it correctly characterises a CNC machining centre and --- on a
single-channel record, having lost the two checks that carry the
discriminating power --- fails to reject a solar inverter, a device that
cannot be shut down at all. A validation battery is only worth quoting
if it can reject something; this one rejects four machines out of eight
and, under reduced instrumentation, is shown to have a blind spot.

\medskip
\noindent
Set against the improvement the framework can extract, the honest
summary is that the facility's total attainable standby saving is under
$90$\,USD/yr, because the plant already de-energises for its long idle
gaps. \textbf{On this facility the framework's value is validated state
inference and decision quality, not recovered energy}, and that
conclusion belongs in the abstract rather than at the end of a
discussion. Whether a plant that does not already shut down manually,
or one with a larger energised-idle fraction, changes the arithmetic is
exactly the question the method now makes answerable from a meter
alone.

\subsection*{Future work}

Four extensions follow directly from the limitations of
Section~\ref{sec:limitations}. A \emph{longer record} is the only thing
that can settle the economic question: eighteen factory days is the
binding constraint on every currency interval here, and no estimator
narrows it. A \emph{second six-channel industrial dataset} would test
the physics battery rather than only the pipeline, which the available
single-channel record cannot do. A \emph{hidden semi-Markov
formulation}~\cite{yu2010hsmm} would learn the state-duration
distribution instead of imposing a constant, and a per-machine state
count selected under the physics veto would remove the $k=4$ assumption
that the eight-machine study shows to be wrong in both directions.
Finally, a \emph{seed-resampled ablation of the validation battery
itself}, check by check, would price the component this paper argues is
most important and is the one component its ablation table does not
contain.

%  Bibliography: paper/references.bib
% =====================================================================

\section{Conclusion}
\label{sec:conclusion}

This paper set out to connect two capabilities that had matured
separately: inferring an industrial machine's operating state from its
electrical measurements, and deciding optimally when to de-energise it.
The connecting quantity is the duration of a non-productive interval,
which the first literature does not forecast and the second assumes is
given. We proposed a unified label-free framework that supplies it, and
evaluated every stage end to end on a public industrial benchmark of
eight machines.

Four results stand.

\textbf{Label-free state inference can be made trustworthy, but only if
the fit is gated.} Fitted at ten seeds on identical data, the labeller
places the standby cluster on a physically different state in three of
them, and the headline standby total moves by a third of its value.
Five circuit-level checks --- none of which uses the model's own
statistics --- separate those three from the other seven exactly, and
conditioning on them collapses the spread to $3\%$. An accept/reject
step is therefore part of the method, not an appendix to it, and this is
the paper's central methodological claim.

\textbf{Temporal coherence requires a duration model, not a stronger
prior.} A transition matrix can penalise a state change by at most
$\log(p_{\text{self}}/p_{\text{switch}})$ --- $9.23$ nats here --- while
the emission term routinely exceeds that by orders of magnitude, so at
over half of all timesteps the transition model cannot influence the
decode at all. Viterbi decoding takes flicker from $66.9\%$ to $44.5\%$;
an explicit minimum-dwell constraint takes it to zero, at a cost of
$0.19\%$ of rows and no measurable change to any physics or schedule
check, on all eight machines.

\textbf{Two of the framework's own components did not earn their
place, and are reported as such.} An untrained persistence reference
beats all five neural sequence forecasters, the originally proposed
sequence-to-sequence model by the largest margin of the seven; only a
gradient-boosted tree clears the reference, and by $1.45$\,s of
per-window MAE. On the economics, the proposed policy cannot be shown to
beat the plant's own behaviour on an 18-day held-out sample. What it
does beat, significantly and in every scenario, is the fixed-threshold
rule it was built to replace. The defensible decision-layer claim is
that \emph{adaptive thresholds beat a fixed one}, not that forecasting
beats not forecasting.

\textbf{The framework says where it does not apply.} Across the plant
the state definition is physically valid on the four motor-driven
machines and is rejected by the battery on the four auxiliary ones,
where no energised-idle state exists to find. Applied to a second site
it correctly characterises a CNC machining centre and --- on a
single-channel record, having lost the two checks that carry the
discriminating power --- fails to reject a solar inverter, a device that
cannot be shut down at all. A validation battery is only worth quoting
if it can reject something; this one rejects four machines out of eight
and, under reduced instrumentation, is shown to have a blind spot.

\medskip
\noindent
Set against the improvement the framework can extract, the honest
summary is that the facility's total attainable standby saving is under
$90$\,USD/yr, because the plant already de-energises for its long idle
gaps. \textbf{On this facility the framework's value is validated state
inference and decision quality, not recovered energy}, and that
conclusion belongs in the abstract rather than at the end of a
discussion. Whether a plant that does not already shut down manually,
or one with a larger energised-idle fraction, changes the arithmetic is
exactly the question the method now makes answerable from a meter
alone.

\subsection*{Future work}

Four extensions follow directly from the limitations of
Section~\ref{sec:limitations}. A \emph{longer record} is the only thing
that can settle the economic question: eighteen factory days is the
binding constraint on every currency interval here, and no estimator
narrows it. A \emph{second six-channel industrial dataset} would test
the physics battery rather than only the pipeline, which the available
single-channel record cannot do. A \emph{hidden semi-Markov
formulation}~\cite{yu2010hsmm} would learn the state-duration
distribution instead of imposing a constant, and a per-machine state
count selected under the physics veto would remove the $k=4$ assumption
that the eight-machine study shows to be wrong in both directions.
Finally, a \emph{seed-resampled ablation of the validation battery
itself}, check by check, would price the component this paper argues is
most important and is the one component its ablation table does not
contain.
